\documentclass[11pt]{article}

\usepackage[margin=1in]{geometry}
\usepackage{amsmath,amssymb,amsthm}
\usepackage{mathtools}
\usepackage{enumitem}
\usepackage{hyperref}

\newtheorem{theorem}{Theorem}[section]
\newtheorem{lemma}[theorem]{Lemma}
\newtheorem{proposition}[theorem]{Proposition}
\newtheorem{corollary}[theorem]{Corollary}
\theoremstyle{definition}
\newtheorem{definition}[theorem]{Definition}
\newtheorem{example}[theorem]{Example}
\theoremstyle{remark}
\newtheorem{remark}[theorem]{Remark}

\newcommand{\Con}{\operatorname{Con}}
\newcommand{\End}{\operatorname{End}}
\newcommand{\rank}{\operatorname{rank}}
\newcommand{\spn}{\operatorname{span}}
\newcommand{\Obs}{O}
\newcommand{\R}{\mathbb{R}}
\newcommand{\eps}{\varepsilon}
\newcommand{\join}{\vee}
\newcommand{\meet}{\wedge}
\newcommand{\J}{\mathsf{J}}
\newcommand{\id}{\mathrm{id}}

\title{Congruence obstructions in exact predictive agents}
\author{}
\date{}

\begin{document}
\maketitle

\begin{abstract}
Canonical predictive states and state-level refinement are standard in
computational mechanics, while finite-rank Hankel minimisation is classical in
weighted-automata and observable-operator-model theory.  This paper lifts an
exact state intertwiner to a surjection between transition monoids.  Congruence
correspondence then identifies the canonical predictive congruence lattice with
an interval in the agent's congruence lattice, so every canonical $M_3$ or
$N_5$ witness is inherited by every reachable exact predictor.

The interpretation of this invariant has a sharp limit.  For any first-order
Markov chain in the stated full-support, distinct-row class, the canonical
transition monoid is a reset monoid.  For three or more symbols its congruence
lattice is non-distributive, although the predictor merely stores the latest
symbol.  Non-distributivity is therefore not by itself evidence of entangled or
uninterpretable internal features.  For controllers, a further corollary is
restricted to the zero-regret, complete-test limit: predictive memory then
refines predictive state, and deterministic counterfactual memory updates
supply the state quotient needed for monoid descent.  Finite regret gives only
a quantitative no-aliasing guarantee, not an exact congruence interval.
\end{abstract}

\section{Question and scope}

Let a world produce a sequence of observations from a finite alphabet $\Obs$.
An exact predictive agent has enough internal information to return the
conditional distribution of every finite future.  Such an agent may nevertheless
have redundant states or implementation-specific dynamics.  Our question is:
\begin{quote}
Which transition structure must every exact predictive agent contain?
\end{quote}

The answer below is a quotient theorem.  The process has a canonical predictive
state machine, and every reachable exact predictor maps onto it.  This yields an
invariant ``world-required'' interval inside the predictor's congruence lattice.
The unconditional theorem concerns prediction, not control.  Section
\ref{sec:controllers} proves that fixed-task optimality is insufficient and
states the additional behavioural-recovery and recurrent-state assumptions
under which a controller inherits the same quotient.  No claim is made that
every competent controller is predictively sufficient or that a neural
network's raw hidden state is automatically a linear realisation.

\section{Prior work and contribution}
\label{sec:positioning}

The construction below combines several established theories:
\begin{enumerate}[label=(\roman*)]
  \item identifying histories by their complete conditional future laws, and
        the fact that any equally predictive rival refines this causal-state
        partition, belong to computational mechanics
        \cite{CrutchfieldYoung1989,ShaliziCrutchfield2001};
  \item the canonical predictive machine for an action--output process is the
        $\varepsilon$-transducer construction
        \cite{BarnettCrutchfield2015};
  \item finite-rank Hankel minimisation and the quotient of any reachable linear
        realisation onto the minimal one are classical weighted-automata,
        rational-series, and observable-operator-model results
        \cite{Fliess1974,Jaeger2000};
  \item behavioural recovery of a controlled transition law from sufficiently
        general goal-conditioned policies is due to Richens et al., with
        stochastic and partially observable extensions due to Cifuentes
        \cite{Richens2025,Cifuentes2026}; transition-local certification is
        developed by Lu et al.\ \cite{LuEtAl2026};
  \item Nayebi establishes quantitative recovery of predictive states and
        no-aliasing constraints on memory under partial observability
        \cite{Nayebi2026}.
\end{enumerate}
Accordingly, the distinction in Section~\ref{sec:controllers} between
behavioural recovery and internal predictive memory is not claimed here as a
new distinction.

The contributions of this paper are narrower and begin at the
transition-monoid layer:
\begin{enumerate}[label=(\arabic*)]
  \item lifting a surjective state intertwiner to a surjection of generated
        transition monoids;
  \item interpreting congruence correspondence as a process-forced interval in
        every exact predictor's congruence lattice;
  \item framing $M_3$ and $N_5$ witnesses in that interval as inherited
        obstructions;
  \item proving the reset-monoid genericity result for the stated full-support
        Markov class, and using it to show that non-distributivity alone is not
        an entanglement certificate;
\end{enumerate}

\section{Stochastic languages and conditional futures}
\label{sec:process}

\begin{definition}[Stochastic language]
A stochastic language over a finite alphabet $\Obs$ is a function
\[
  p:\Obs^*\longrightarrow[0,1]
\]
such that
\[
  p(\eps)=1,
  \qquad
  p(u)=\sum_{o\in\Obs}p(uo)
  \quad (u\in\Obs^*).
\]
Here $p(u)$ is the probability that the process begins with the word $u$.
\end{definition}

Right consistency is all that is needed below.  If the process is also stationary,
then its cylinder probabilities additionally satisfy
\[
  p(u)=\sum_{o\in\Obs}p(ou).
\]
Stationarity may be imposed as a restriction on the class of processes, but none
of the quotient arguments uses it.

For a history $u$ with $p(u)>0$, define its conditional future law
\[
  q_u(v)=\frac{p(uv)}{p(u)}
  \qquad(v\in\Obs^*).
\]
Then $q_u(\eps)=1$ and $q_u(v)=\sum_o q_u(vo)$.  Histories $u,u'$ are
predictively equivalent precisely when $q_u=q_{u'}$.

If $p(u)=0$, right consistency and non-negativity imply $p(uv)=0$ for every
$v$.  Conditional prediction after $u$ is undefined.  We represent every such
impossible continuation by one absorbing formal symbol
$\dagger\notin X_p$.

\begin{definition}[Canonical predictive states]
The predictive state set of $p$ is
\[
  X_p=\{q_u:p(u)>0\}.
\]
Write $\widehat X_p=X_p\cup\{\dagger\}$.
\end{definition}

This is the usual predictive or causal-state identification by complete future
law \cite{CrutchfieldYoung1989,ShaliziCrutchfield2001}.

For $o\in\Obs$, define $\tau_o:\widehat X_p\to\widehat X_p$ by
\[
  \tau_o(q_u)=
  \begin{cases}
    q_{uo},&p(uo)>0,\\
    \dagger,&p(uo)=0,
  \end{cases}
  \qquad
  \tau_o(\dagger)=\dagger.
\]

\begin{lemma}
The map $\tau_o$ is well-defined: if $q_u=q_{u'}$, then either both
$p(uo)$ and $p(u'o)$ vanish, or $q_{uo}=q_{u'o}$.
\end{lemma}

\begin{proof}
The one-symbol probabilities agree:
\[
  \frac{p(uo)}{p(u)}=q_u(o)=q_{u'}(o)
  =\frac{p(u'o)}{p(u')}.
\]
Thus one is zero exactly when the other is.  If they are positive, then for every
$v$,
\[
  q_{uo}(v)=\frac{q_u(ov)}{q_u(o)}
  =\frac{q_{u'}(ov)}{q_{u'}(o)}
  =q_{u'o}(v).
\]
\end{proof}

\begin{definition}[Canonical predictive transition monoid]
Let
\[
  S_p=\langle\tau_o:o\in\Obs\rangle
  \subseteq\mathcal T_{\widehat X_p}
\]
be the transformation monoid generated by the predictive updates.  Composition is
ordinary function composition.  For a word $w=o_1\cdots o_n$, write
$\tau_w=\tau_{o_n}\circ\cdots\circ\tau_{o_1}$, so that
$\tau_w(q_u)=q_{uw}$ whenever $p(uw)>0$.
\end{definition}

\section{Every exact predictive agent contains the canonical machine}
\label{sec:agent}

We model only the part of an agent used for exact prediction.

\begin{definition}[Reachable exact predictive agent]
A reachable exact predictive agent for $p$ consists of
\[
  A=(Q,q_0,d,\{\delta_o\}_{o\in\Obs},F).
\]
For $u=o_1\cdots o_n$, write
\[
  \delta_u=\delta_{o_n}\circ\cdots\circ\delta_{o_1},
  \qquad \delta_\eps=\id.
\]
The non-dead reachable states are
\[
  Q^+=\{\delta_u(q_0):p(u)>0\},
\]
the dead state $d$ is not in $Q^+$, and
$Q=Q^+\mathbin{\dot\cup}\{d\}$.  Each $\delta_o:Q\to Q$ satisfies
$\delta_o(d)=d$ and
\[
  p(u)>0,\ p(uo)=0
  \quad\Longrightarrow\quad
  \delta_o(\delta_u(q_0))=d.
\]
Finally, $F:Q^+\to X_p$ returns the state's complete conditional future law.
The predictor is exact when, for every $u$ with $p(u)>0$,
\[
  F(\delta_u(q_0))=q_u.
\]
\end{definition}

Trimming impossible branches changes no prediction made on the process.  It makes
the update maps total and avoids assigning arbitrary semantics to conditionals on
null events.

\begin{theorem}[Canonical predictor quotient]
\label{thm:predictor-quotient}
Let $A$ be a reachable exact predictive agent for $p$.  There is a surjection
\[
  L:Q\twoheadrightarrow\widehat X_p
\]
such that
\[
  L\circ\delta_o=\tau_o\circ L
  \qquad(o\in\Obs).
\]
\end{theorem}

\begin{proof}
For a non-dead reachable state $q=\delta_u(q_0)$, set $L(q)=q_u$, and send the
agent's dead state to $\dagger$.  This is well-defined: if
$\delta_u(q_0)=\delta_v(q_0)$, exactness gives
$q_u=F(\delta_u(q_0))=F(\delta_v(q_0))=q_v$.

Every element $q_u\in X_p$ is hit by the state $\delta_u(q_0)$, so $L$ is
surjective.  If $p(uo)>0$, then
\[
  L(\delta_o\delta_u(q_0))=q_{uo}
  =\tau_o(q_u)=\tau_oL(\delta_u(q_0)).
\]
If $p(uo)=0$, both sides are $\dagger$ by the trimming convention.  The equality
also holds at the absorbing dead state.
\end{proof}

At the level of reachable history partitions, this is the pointwise analogue
of the Refinement Lemma: any prescient rival partition refines the causal-state
partition, up to null histories \cite{ShaliziCrutchfield2001}.  The next step
lifts that state-level refinement to transition monoids.

Let $S_A=\langle\delta_o:o\in\Obs\rangle\subseteq\mathcal T_Q$.

\begin{theorem}[Transition-monoid quotient]
\label{thm:monoid-quotient}
Under the hypotheses of Theorem~\ref{thm:predictor-quotient}, the assignment
\[
  \Phi(\delta_w)=\tau_w
\]
defines a surjective monoid homomorphism
\[
  \Phi:S_A\twoheadrightarrow S_p.
\]
\end{theorem}

\begin{proof}
Iterating the intertwining equations gives
\[
  L\circ\delta_w=\tau_w\circ L
  \qquad(w\in\Obs^*).
\]
If $\delta_w=\delta_v$, then
\[
  \tau_w\circ L=L\circ\delta_w=L\circ\delta_v=\tau_v\circ L.
\]
Surjectivity of $L$ implies $\tau_w=\tau_v$, so $\Phi$ is well-defined.  It
also respects composition: for words $v,w$ our convention gives
\[
  \delta_w\circ\delta_v=\delta_{vw},
  \qquad
  \tau_w\circ\tau_v=\tau_{vw},
\]
and hence
\[
  \Phi(\delta_w\circ\delta_v)
  =\tau_{vw}
  =\Phi(\delta_w)\circ\Phi(\delta_v).
\]
It is surjective because it hits every generator $\tau_o$.
\end{proof}

The abstract descent step is also formalised in Lean.\footnote{\url{https://github.com/ayurpulle/DovetailProofs/blob/main/DovetailProofs/Descent.lean}}
That development uses \texttt{FreeMonoid.lift}, whose endomorphism convention is
opposite to the temporal convention above.  It therefore constructs the
opposite transition monoid; the congruence conclusions transfer because
$\Con(S)=\Con(S^{\mathrm{op}})$.

\begin{corollary}[Congruence inheritance]
\label{cor:con-inheritance}
Let $\kappa=\ker\Phi$.  Then
\[
  \Con(S_p)\cong[\kappa,\top]\subseteq\Con(S_A).
\]
Consequently, if $\Con(S_p)$ is non-distributive, then $\Con(S_A)$ is
non-distributive.  More specifically, every $M_3$ or $N_5$ sublattice of
$\Con(S_p)$ occurs inside the interval $[\kappa,\top]$ of the agent's
congruence lattice.
\end{corollary}

\begin{proof}
The first statement is the congruence correspondence theorem for the surjection
$\Phi$.  An interval is closed under lattice meets and joins, so any
non-distributive witness in the interval is also a witness in the whole lattice.
\end{proof}

This is the agent-facing result: redundant exact predictors may add
implementation-specific distinctions below $\kappa$, but cannot remove the
canonical predictive interval above it.

\section{The finite-rank linear realisation}
\label{sec:linear}

The preceding theorem does not assume that the agent is linear.  The Hankel
construction gives a canonical finite-dimensional realisation when $p$ has finite
Hankel rank \cite{Fliess1974,Jaeger2000}.

Let
\[
  W=\R^{\Obs^*}
\]
be the vector space of all real functions on finite futures.  We do not choose a
Hamel basis for $W$; the coordinate delta-functions span only the finitely
supported functions.  Define the unnormalised Hankel row
\[
  r_u(v)=p(uv)
\]
and the linear realisation space
\[
  V=\spn\{r_u:u\in\Obs^*\}\subseteq W.
\]
For the infinite Hankel matrix $H_{u,v}=p(uv)$, define
\[
  \rank H:=\dim\spn\{r_u:u\in\Obs^*\}.
\]
Assume in this section that $\rank H<\infty$.  Then
$\dim V=\rank H$.
The rows $r_u$ are not conditional predictive states: for $p(u)>0$,
\[
  r_u=p(u)q_u.
\]
The assignment
\[
  \beta:X_p\longrightarrow
  \{\text{nonzero reachable positive Hankel rays}\},
  \qquad
  \beta(q_u)=[r_u]=\{c r_u:c>0\},
\]
is a well-defined bijection.  Indeed, proportional rows become equal after
normalising their $\eps$-coordinate to one, which recovers $q_u$.  Moreover,
$\beta$ intertwines the conditional and projective updates:
\[
  \beta(\tau_o(q_u))=T_o\cdot\beta(q_u)
\]
whenever $p(uo)>0$, with impossible extensions sent to $\dagger$ on both
sides.  We henceforth identify $X_p$ with this ray set.  The vector space $V$,
by contrast, is the minimal linear realisation space.

For $o\in\Obs$, define
\[
  (T_of)(v)=f(ov).
\]
Then $T_or_u=r_{uo}$, so $T_o$ restricts to a linear endomorphism of $V$.
We use the column convention.  For $w=o_1\cdots o_n$, put
\[
  T_w=T_{o_n}\cdots T_{o_1};
\]
then $T_wr_u=r_{uw}$.  Hence $w\mapsto T_w$ is an anti-homomorphism into
$\End(V)$ under ordinary matrix multiplication.  Equivalently it is a
homomorphism into the opposite monoid.  This choice has no effect on congruence
lattices:
\[
  \Con(S)=\Con(S^{\mathrm{op}}),
\]
because the same equivalence relations are compatible with multiplication in
both orders.

The readout is the linear functional
\[
  \lambda:V\to\R,\qquad\lambda(f)=f(\eps),
\]
and
\[
  p(w)=\lambda(T_wr_\eps).
\]
In a general basis $\lambda$ is represented by a basis-dependent covector, not
necessarily by an all-ones vector.

\begin{proposition}
The realisation $(V,r_\eps,\{T_o\},\lambda)$ is reachable and observable, and
therefore minimal.  Every reachable linear realisation of $p$ has dimension at
least $\dim V$.
\end{proposition}

\begin{proof}
Reachability is the definition $V=\spn\{T_ur_\eps:u\in\Obs^*\}$.  If
$f\in V$ satisfies $\lambda(T_vf)=0$ for every $v$, then
$f(v)=0$ for every future $v$, so $f=0$; this is observability.  The
surjection constructed in Proposition~\ref{prop:linear-quotient} below gives
the dimension bound.
\end{proof}

\begin{proposition}[Linear predictor quotient]
\label{prop:linear-quotient}
Let $V'$ be a vector space, $x_0\in V'$, $A_o\in\End(V')$, and
$\lambda'\in(V')^*$.  For $w=o_1\cdots o_n$, set
\[
  A_w=A_{o_n}\cdots A_{o_1}.
\]
Suppose these data are a reachable linear realisation of the same stochastic
language:
\[
  p(w)=\lambda'(A_wx_0)
  \quad(w\in\Obs^*),
  \qquad
  V'=\spn\{A_ux_0:u\in\Obs^*\}.
\]
Then there is a surjective linear map $L:V'\twoheadrightarrow V$ satisfying
\[
  LA_o=T_oL
  \qquad(o\in\Obs).
\]
It therefore induces a surjective homomorphism from the agent's operator monoid
onto the minimal operator monoid.
\end{proposition}

\begin{proof}
For $x\in V'$, define its future behaviour by
\[
  (Lx)(v)=\lambda'(A_vx).
\]
Reachability writes every $x$ as a linear combination of states $A_ux_0$, and
\[
  L(A_ux_0)=r_u.
\]
Thus $Lx\in V$ and $L$ is onto.  Moreover,
\[
  (LA_ox)(v)=\lambda'(A_vA_ox)
  =(T_oLx)(v),
\]
so $LA_o=T_oL$.  The induced monoid map is well-defined by the same descent
argument as in Theorem~\ref{thm:monoid-quotient}; explicitly,
\[
  \Psi:\langle A_o:o\in\Obs\rangle\twoheadrightarrow
  \langle T_o:o\in\Obs\rangle,
  \qquad
  \Psi(A_w)=T_w.
\]
\end{proof}

This is the linear analogue of the predictor quotient at the level of
unnormalised joint-future behaviour.  Normalising positive reachable rows and
collapsing null histories recovers the conditional-state quotient.  The result
applies to exact weighted automata, observable-operator models, and other
reachable linear predictors of the complete joint language.

\section{Conditional-state dynamics versus quantitative prediction}
\label{sec:quantitative}

Let $S_{\mathrm{lin}}=\langle T_o:o\in\Obs\rangle\subseteq\End(V)$.  Its action
on the positive Hankel rays is
\[
  T\cdot[r_u]=
  \begin{cases}
    [Tr_u],&Tr_u\ne0,\\
    \dagger,&Tr_u=0,
  \end{cases}
  \qquad T\cdot\dagger=\dagger.
\]
This gives a monoid homomorphism
\[
  \alpha:S_{\mathrm{lin}}\longrightarrow
  \mathcal T_{\widehat X_p}.
\]
Its image is the canonical predictive transition monoid $S_p$ and
\[
  \kappa_{\mathrm{pred}}=\ker\alpha
\]
identifies operators that induce the same conditional-state update, including
the same pattern of impossible continuations.

There is a simple computable lower bound.  Define
\[
  T\mathrel{\theta_{\mathrm{scal}}}T'
  \quad\Longleftrightarrow\quad
  T'=cT\text{ for some }c>0.
\]
This is a congruence on $S_{\mathrm{lin}}$: positive scalar multiplication is
compatible with composition on both sides.  Moreover,
\[
  \theta_{\mathrm{scal}}\subseteq\kappa_{\mathrm{pred}}.
\]
Equality need not hold.  Agreement on a spanning set of rays does not force
global proportionality: as a linear-algebra illustration, $I$ and
$\operatorname{diag}(2,3)$ agree on $[e_1]$ and $[e_2]$ but are not scalar
multiples.  Section~\ref{sec:reset} gives a process-realizable strict inclusion.

The quotient by $\kappa_{\mathrm{pred}}$ can forget quantitative information.  For
$x=[r_u]\in X_p$, define the likelihood cocycle
\[
  \ell_T(x)=\frac{\lambda(Tr_u)}{\lambda(r_u)},
\]
with value $0$ when $Tr_u=0$, and set $\ell_T(\dagger)=1$.  This is independent
of the positive representative of the ray and satisfies
\[
  \ell_{TU}(x)=\ell_T(U\cdot x)\ell_U(x).
\]
If $T=T_w$, then
\[
  \ell_{T_w}([r_u])=\frac{p(uw)}{p(u)},
\]
the probability of observing $w$ from predictive state $q_u$.

\begin{proposition}[Ray action plus likelihood is faithful]
\label{prop:faithful}
If $T,T'\in S_{\mathrm{lin}}$ have the same action on $\widehat X_p$ and
\[
  \ell_T(x)=\ell_{T'}(x)
  \qquad(x\in X_p),
\]
then $T=T'$.
\end{proposition}

\begin{proof}
Fix $u$ with $p(u)>0$.  If $Tr_u=0$, equality of the ray actions gives
$T'r_u=0$.  Otherwise the two outputs lie on the same positive ray.  Equality of
likelihoods says that they also have the same value under $\lambda$.  A nonzero
Hankel row has positive $\eps$-coordinate, so its positive ray together with that
coordinate determines it uniquely.  Hence $Tr_u=T'r_u$.

The nonzero rows $r_u$ span $V$; zero rows contribute nothing.  Therefore $T$ and
$T'$ agree on a spanning set and are equal.
\end{proof}

Thus $S_p\cong S_{\mathrm{lin}}/\kappa_{\mathrm{pred}}$ is the canonical
\emph{conditional-state transition skeleton}, not the complete stochastic
process.  The concrete minimal realisation, with its distinguished generators,
initial state, and readout, packages both the normalized update and its
likelihood.  For an i.i.d.\ process, $X_p$ has only one predictive state, so
every positive word acts identically on $X_p$, although different Bernoulli
parameters and different words can have different probabilities.

\section{A simple Markov chain already gives non-distributivity}
\label{sec:reset}

The next example is both process-realizable and transparent to an exact agent.
It shows why non-distributivity cannot by itself certify representational
entanglement.

Let $\Obs=\{1,2,3\}$ and consider the stationary first-order Markov chain with
uniform stationary distribution and transition matrix
\[
  M=
  \begin{pmatrix}
    \frac12&\frac13&\frac16\\[2pt]
    \frac13&\frac16&\frac12\\[2pt]
    \frac16&\frac12&\frac13
  \end{pmatrix}.
\]
The matrix has full support, and its three rows are distinct.  After any nonempty
history, the conditional future law is determined by the last observed symbol.
Write $g_i$ for the predictive state after observing $i$.  The initial
stationary predictive state has uniform one-symbol marginal and is distinct from
the three $g_i$, so
\[
  X_p=\{q_\eps,g_1,g_2,g_3\}.
\]

Observing $i$ sends every state in $X_p$ to $g_i$ and fixes $\dagger$.  Let
$c_i$ denote this reset map.  The canonical predictive transition monoid is
\[
  R_3^1=\{1,c_1,c_2,c_3\},
  \qquad c_i c_j=c_i,
\]
where the displayed multiplication uses ordinary composition, with the left
factor applied last.

This example also makes
$\theta_{\mathrm{scal}}\subsetneq\kappa_{\mathrm{pred}}$ concrete.  The
restrictions of $g_1,g_2,g_3$ to one-symbol futures are the rows of $M$, and
$\det M=-1/12$, so these states are linearly independent.  In that basis,
\[
  T_2g_i=M_{i2}g_2,
  \qquad
  T_{12}g_i=M_{i1}M_{12}g_2.
\]
The coefficient vectors are
\[
  \left(\frac13,\frac16,\frac12\right)
  \quad\text{and}\quad
  \left(\frac16,\frac19,\frac1{18}\right),
\]
which are not proportional.  Hence $T_2$ and $T_{12}$ are not scalar
multiples, although both reset every state in $X_p$ to $g_2$ and fix
$\dagger$.

\begin{theorem}[Congruences of the reset monoid]
\label{thm:reset-con}
For $n\ge1$, let
\[
  R_n^1=\{1,c_1,\ldots,c_n\},
  \qquad c_i c_j=c_i,
\]
with $1$ an identity.  Every non-universal congruence has $\{1\}$ as a singleton
class and restricts to an arbitrary equivalence relation on
$\{c_1,\ldots,c_n\}$.  Consequently,
\[
  \Con(R_n^1)\cong \Pi_n\oplus\mathbf 1,
\]
the partition lattice on $n$ points with one new top element adjoined.
\end{theorem}

\begin{proof}
First take any partition of the constants and keep $1$ in a singleton class.
It is a congruence.  Left multiplication of equivalent constants by $1$ preserves
their equivalence, left multiplication by a constant makes both products equal,
and the analogous right-multiplication checks follow from $c_ic_j=c_i$.

Conversely, suppose a congruence identifies $1$ with some $c_i$.  For every $j$,
right multiplication by $c_j$ gives
\[
  c_j=1c_j\equiv c_ic_j=c_i.
\]
Thus all constants are equivalent to $c_i$, and hence to $1$; the congruence is
universal.  Every non-universal congruence therefore fixes $1$, after which its
restriction to the constants is an arbitrary partition.
\end{proof}

\begin{corollary}
For the three-symbol Markov chain above, $\Con(S_p)$ contains
\[
  \Pi_3\cong M_3
\]
as the interval below the congruence that merges all three reset maps while
leaving the identity separate.  In particular, $\Con(S_p)$ is
non-distributive.
\end{corollary}

\begin{proposition}[Full-support Markov chains give reset monoids]
\label{prop:markov-reset-general}
Let $p$ be a first-order Markov chain on $n\ge2$ observation symbols whose
initial distribution has full support and whose transition matrix has full
support and pairwise distinct rows.  Then
\[
  S_p\cong R_n^1.
\]
Consequently,
\[
  \Con(S_p)\cong\Pi_n\oplus\mathbf 1,
\]
and it is non-distributive for every $n\ge3$.
\end{proposition}

\begin{proof}
After every nonempty history, the conditional future law is the law $g_i$
associated with its last symbol $i$.  Distinct transition rows make the $g_i$
pairwise distinct.  Full support of both the initial distribution and the
transition matrix ensures that observing $i$ is possible from every live
predictive state and resets that state to $g_i$.  Thus the generated
transformations are precisely the identity and the $n$ distinct reset maps
$c_i$, which satisfy $c_ic_j=c_i$.  The congruence statement follows from
Theorem~\ref{thm:reset-con}.
\end{proof}

An exact predictor for this process only needs to retain the latest symbol (and
an initial state before the first observation).  Its update rule is a reset, and
its predictions are the corresponding row of $M$ followed by the same Markov
law.  This representation is finite, explicit, and transparent.  Nevertheless
its canonical transition monoid has an $M_3$ congruence interval.

Combining this example with Corollary~\ref{cor:con-inheritance} gives two
conclusions:
\begin{enumerate}[label=(\roman*)]
  \item every reachable exact predictor of this process inherits a
        non-distributive congruence interval;
  \item that non-distributivity cannot, by itself, imply that the predictor uses
        entangled, distributed, or uninterpretable features.
\end{enumerate}
The first conclusion is an invariant theorem.  The second is a limit on its
interpretation.

\section{What distributivity does say}
\label{sec:coordinates}

The coordinate interpretation is precise only after fixing a finite lattice.

\begin{definition}
In a finite lattice $L$, an element $j\ne\bot$ is join-irreducible if
$j=a\join b$ implies $j=a$ or $j=b$.  Equivalently, $j$ has a unique lower
cover.
\end{definition}

Let $\J(L)$ be the poset of join-irreducible elements, and write
\[
  \J(x)=\{j\in\J(L):j\le x\}.
\]
Every element of a finite lattice is the join of the join-irreducibles below it.
Therefore
\[
  x\longmapsto\J(x)
\]
is injective for \emph{every} finite lattice, not only a distributive one.

\begin{theorem}[Finite Birkhoff representation {\cite{DaveyPriestley2002}}]
For a finite lattice $L$, the following are equivalent:
\begin{enumerate}[label=(\roman*)]
  \item $L$ is distributive;
  \item every join-irreducible is join-prime;
  \item the map above is a lattice isomorphism from $L$ to the down-set lattice
        of $\J(L)$.
\end{enumerate}
Equivalently, $L$ is non-distributive exactly when it contains $M_3$ or $N_5$ as
a sublattice.
\end{theorem}

What can fail without distributivity is not bit-vector injectivity but the
union-composition identity
\[
  \J(a\join b)=\J(a)\cup\J(b).
\]
In $M_3$, the five lattice elements still have five distinct bit-vectors, but
joining two atoms activates the third.  Thus the canonical join-irreducible
coordinates are not union-compositional; equivalently, there is no
union/intersection representation of the lattice.

If $L$ is finite and distributive, then $\J(L)$ being an antichain is equivalent
to $L$ being Boolean.  This is a statement about the lattice only.  Even
$\Con(S)\cong2^n$ does not by itself prove that $S$ decomposes into independent
subsystems.  A direct-product conclusion requires suitable complementary factor
congruences and a congruence-permutability or Chinese-remainder condition.

When $S_p$ is infinite, the finite Birkhoff theorem cannot be applied merely by
restricting attention to finite-index congruences: they form a sublattice, but
that sublattice may itself be infinite.  In this draft the coordinate theorem is
used only for explicitly finite monoids or explicitly finite quotient intervals.

\section{Why congruences are the admissible operator lumpings}
\label{sec:congruences}

The interpretation of $\Con(S)$ concerns transformations, not partitions of
predictive states.  We make the modelling criterion explicit.  An admissible
lumping $\sim$ of elements of a transition monoid $S$ must satisfy:
\begin{enumerate}[label=(D\arabic*)]
  \item it is an equivalence relation, so each transformation has one coarse
        description;
  \item it is dynamically closed: $s\sim t$ implies
        $asb\sim atb$ for every $a,b\in S$.
\end{enumerate}
Condition (D2) says that substituting indistinguishable dynamics into any past
and future dynamical context cannot make them distinguishable.  Conditions
(D1)--(D2) are exactly the definition of a two-sided monoid congruence.

This does not prove that every algebraic congruence is physically or cognitively
salient.  It proves the conditional statement: once admissibility is defined as
context-stable lumping of transition operators, the admissible hierarchy is
$\Con(S)$.  Coarse-graining predictive \emph{states} is a different construction
and should not be conflated with it.

\section{Optimal controllers and the predictive quotient}
\label{sec:controllers}

Actions require a controlled version of the predictive machine.  Fix finite
action and output alphabets $A$ and $O$.  An anticipation-free interface
$\mathcal I$ specifies probabilities
\[
  \mathcal I(o_{1:n}\mid a_{1:n})
\]
which are normalized and causally prefix-consistent:
\[
  \mathcal I(\eps\mid\eps)=1,
  \qquad
  \sum_{o_n\in O}\mathcal I(o_{1:n}\mid a_{1:n})
  =\mathcal I(o_{1:n-1}\mid a_{1:n-1})
  \quad(n\ge1).
\]
A history
\[
  h=(a_1,o_1)\cdots(a_n,o_n)
\]
is admissible when $\mathcal I(o_{1:n}\mid a_{1:n})>0$.  Its residual
interface is
\[
  \mathcal I^h(v\mid u)
  =
  \frac{\mathcal I(o_{1:n}v\mid a_{1:n}u)}
       {\mathcal I(o_{1:n}\mid a_{1:n})},
\]
for equally long future action and output words $u$ and $v$.  Define
\[
  X_{\mathcal I}=\{\mathcal I^h:h\text{ admissible}\}.
\]
Let
\[
  \widehat X_{\mathcal I}
  =X_{\mathcal I}\mathbin{\dot\cup}\{\dagger\}.
\]
Appending an action-output pair gives the total canonical update
\[
  \tau_{a,o}(\mathcal I^h)=
  \begin{cases}
    \mathcal I^{h(a,o)},&\mathcal I^h(o\mid a)>0,\\
    \dagger,&\mathcal I^h(o\mid a)=0,
  \end{cases}
  \qquad
  \tau_{a,o}(\dagger)=\dagger.
\]
This is well-defined: equal residual interfaces have equal one-step output
probabilities, and Bayes conditioning gives equal successor residuals whenever
that probability is positive.  The controlled predictive transition monoid is
\[
  S_{\mathcal I}=\langle\tau_{a,o}:(a,o)\in A\times O\rangle.
\]
This is the $\varepsilon$-transducer construction for controlled processes
\cite{BarnettCrutchfield2015}.

\subsection{Two meanings of ``the controller contains a model''}

The following distinction organises the implications needed for monoid descent;
it is not claimed as new.
\begin{enumerate}[label=(\roman*)]
  \item \emph{Behavioural containment:} the environment's transition law or
        residual interface can be recovered from the controller's complete
        goal-conditioned policy.
  \item \emph{Internal quotient containment:} there is a surjective map from
        reachable controller states to $X_{\mathcal I}$ which commutes with
        every action-output update.
\end{enumerate}
The second statement implies a structural quotient of transition monoids.  The
first is an identifiability statement about behaviour and does not, by itself,
locate a model in the controller's activations or recurrent state.
For behavioural containment, the recovery functional must be fixed uniformly
over a class of environments sharing the same action, output, goal, and reward
semantics; allowing it to depend on the unknown true interface would make the
claim vacuous.

\subsection{Why optimality alone is insufficient}

The proof uses the controller conventions of
Subsection~\ref{subsec:controller-formalism}.  States are reachable, the dead
state is absorbing, and updates are total on counterfactual histories.

\begin{proposition}[Fixed-task optimality is insufficient]
\label{prop:controller-no-go}
Optimality for a single reward function does not force either form of
containment.
\end{proposition}

\begin{proof}
Take a full-support controlled process with at least two distinct predictive
states and give every trajectory reward zero.  Every policy is optimal,
including a controller with one live state (and the formal dead state) which
chooses one fixed action after every counterfactual history.  A dead-state
preserving map from this controller cannot surject onto two or more live
predictive states.  Moreover, the same constant policy is optimal for a family
of full-support environments with common action and output alphabets but
different transition laws.  No recovery functional fixed across that family
can recover the environment from the policy.
\end{proof}

The counterexample is not merely an artifact of a constant reward.  Histories
with different future laws may be control-equivalent for a particular task, and
an optimal controller may merge them.  Richens et al.\ prove a stronger
policy-level negative result for the class
$\Psi_{\mathrm{myopic}}$ of depth-one goals requiring attainment immediately
after the first action: even a policy optimal for every such goal identifies no
nontrivial entrywise bound on the environment's transition probabilities; the
trivial error bound $1$ is tight \cite{Richens2025}.  Multi-step task diversity,
not optimality alone, is the essential extra condition.

\subsection{Behavioural recovery: Richens and Cifuentes}

Richens et al.\ study a fixed finite, communicating, stationary controlled
Markov process with at least two actions, assumed fully observed by the agent.
Their policies are deterministic, history-dependent, and goal-conditioned.
For a goal $\psi$, their relative-failure condition is
\[
  \Pr(\psi\mid\pi,s_0)
  \ge(1-\delta)\sup_{\pi'}\Pr(\psi\mid\pi',s_0).
\]
Their main recovery theorem can be summarised as follows.

\begin{theorem}[Policy-level world-model recovery {\cite{Richens2025}}]
\label{thm:richens}
For the class above, suppose the relative-failure condition holds for every
composite goal in $\Psi_n$, the goals of depth at most $n>1$, with
$0\le\delta<1$.  The policy determines an estimate $\widehat P$ of the
transition kernel satisfying
\[
  \left|\widehat P_{ss'}(a)-P_{ss'}(a)\right|
  \le
  \sqrt{\frac{2P_{ss'}(a)(1-P_{ss'}(a))}
  {(n-1)(1-\delta)}}.
\]
Thus the guaranteed error decreases as $\delta$ decreases or $n$ increases,
and tends to zero for fixed $\delta<1$ as $n\to\infty$.
\end{theorem}

This is policy-level, bounded-error recovery.  At finite goal depth the theorem
guarantees an approximation, including for an optimal policy; it does not rule
out exact identification in special cases.  For fixed $\delta<1$, it guarantees
arbitrarily accurate recovery as the goal depth grows.  The recovery map is
from the policy, not from the policy's implementation, so a further bridge is
needed for an internal quotient.

Cifuentes removes two restrictions from this behavioural result.  The queried
goal-conditioned policy may be stochastic, and it may act from observation
histories in a partially observable controlled Markov process
\cite{Cifuentes2026}.  The result is still black-box recovery: suitable
goal-conditioned policy queries determine an approximation to the latent
transition kernel.  It does not identify where that information is stored in a
particular recurrent implementation.

There is also an important scope restriction.  In Cifuentes's partially
observable theorem, goals are still formulated and evaluated using the latent
environment states, even though the policy observes only outputs.  Cifuentes
notes that if goals themselves are restricted to observation-level events, the
latent transition kernel is not generally identifiable and leaves the
recoverable part of the environment open.  Thus this theorem removes the
full-observability assumption on the \emph{policy}, but it does not by itself
recover an arbitrary residual action--output interface from purely
observation-grounded goals.

\subsection{Internal predictive-memory necessity: Nayebi}

Nayebi takes a representation-level route rather than first recovering a latent
transition table \cite{Nayebi2026}.  In a POMDP, a predictive test
$T=(\alpha,W)$ prescribes a future action word $\alpha$ and asks whether the
resulting output word lies in an event $W$.  Its success probability after
history $h$ is
\[
  p_T(h)=\mathcal I^h(W\mid\alpha).
\]
For a test family $\mathcal T$, write
\[
  \eta_{\mathcal T}(h)
  =
  \bigl(p_T(h)\bigr)_{T\in\mathcal T}.
\]
Threshold betting goals compare $p_T(h)$ with a known lottery
$\lambda\in[0,1]$.  The three results directly relevant to the bridge are:
\begin{enumerate}[label=(\roman*)]
  \item Nayebi's Theorem 3 shows that low average regret over a grid of threshold bets
        quantitatively recovers the tested predictive coordinates
        $\eta_{\mathcal T}(h)$;
  \item Nayebi's Theorem 4 shows that, under a finite-dimensional linear-PSR hypothesis
        and a full-rank choice of histories, the same method quantitatively
        recovers the linear predictive-state operators;
  \item Nayebi's Theorem 5 gives the quantitative no-aliasing bound for any memory
        statistic through which the policy factors.
\end{enumerate}

More directly for internal memory, suppose a policy factors through a memory
statistic $M(h)$.  If $M(h)=M(h')$ while some test requires opposite bets at
the two histories with margin at least $\gamma$, Nayebi's no-aliasing theorem
gives
\[
  \overline\delta_{\mathcal P}(\pi)
  \ge
  q^{\mathrm{Alias}}_\gamma(M)\,
  \frac{c(\gamma)}2,
  \qquad
  c(\gamma)=\frac{4\gamma}{1+2\gamma},
\]
where $q^{\mathrm{Alias}}_\gamma(M)$ is the evaluated mass of such witnessed
aliases.  Hence low regret quantitatively limits predictive aliasing.  In the
zero-regret limit this yields an exact refinement, but only under explicit
completeness assumptions.

\begin{corollary}[Zero-regret predictive-memory refinement]
\label{cor:nayebi-refinement}
Let $H_0$ be a set of evaluated histories and let a goal-conditioned policy
factor through a memory statistic $M:H_0\to Q^+$.  Assume the memory carrier is
reachable:
\[
  Q^+=\{M(h):h\in H_0\},
  \qquad
  Q=Q^+\mathbin{\dot\cup}\{d\}.
\]
Suppose:
\begin{enumerate}[label=(\roman*)]
  \item at every $h\in H_0$, the policy has zero regret on every available
        threshold bet;
  \item the available threshold bets separate unequal predictive-coordinate
        vectors: whenever
        $\eta_{\mathcal T}(h)\ne\eta_{\mathcal T}(h')$, some
        $T\in\mathcal T$ and $\lambda\in[0,1]$ satisfy either
        \[
          p_T(h)>\lambda>p_T(h')
        \]
        or the reverse inequality.
\end{enumerate}
Then, for all $h,h'\in H_0$,
\[
  M(h)=M(h')
  \quad\Longrightarrow\quad
  \eta_{\mathcal T}(h)=\eta_{\mathcal T}(h').
\]
If $\mathcal T$ is predictively complete, so equality of its coordinate
vectors is equivalent to equality of residual interfaces, then
\[
  L(M(h))=\mathcal I^h,
  \qquad
  L(d)=\dagger
\]
defines a well-defined surjection
\[
  L:Q\twoheadrightarrow
  \{\mathcal I^h:h\in H_0\}\mathbin{\dot\cup}\{\dagger\}.
\]

If in addition $H_0$ is the full admissible support (or is transition-closed
after adjoining a dead state), and there are deterministic, goal-independent
updates $\delta_{a,o}:Q\to Q$ defined for every counterfactual
action--output pair such that
\[
  \delta_{a,o}(M(h))=M(h(a,o))
\]
for admissible extensions, while impossible extensions go to an absorbing
state $d$, then
\[
  L\circ\delta_{a,o}=\tau_{a,o}\circ L.
\]
For the full support this is the state quotient required for monoid descent,
and hence it induces the monoid surjection and congruence interval of
Corollary~\ref{cor:controller-con}.
\end{corollary}

\begin{proof}
Assume $M(h)=M(h')$ but
$\eta_{\mathcal T}(h)\ne\eta_{\mathcal T}(h')$.  By separation, choose
$T,\lambda$ with, say,
$p_T(h)>\lambda>p_T(h')$.  The unique optimal threshold bet at $h$ chooses
the test, while the unique optimal bet at $h'$ chooses the lottery.  Because
the policy factors through $M$, equal memory states give the same distribution
over these two choices.  It therefore cannot have zero regret at both
histories, a contradiction.  This proves the displayed refinement.

Predictive completeness makes $L$ well-defined.  The reachability clause
defines it on every element of $Q^+$, and $d$ is assigned separately.  Every
evaluated predictive state has a history representative, so $L$ is
surjective.  For an admissible extension $h(a,o)$, deterministic recurrence gives
\[
  L(\delta_{a,o}M(h))
  =L(M(h(a,o)))
  =\mathcal I^{h(a,o)}
  =\tau_{a,o}L(M(h)).
\]
The transition-closure and dead-state assumptions cover every update.  The
monoid and congruence conclusions then follow from the existing descent and
correspondence results.
\end{proof}

\begin{remark}[Finite regret]
Theorem 5 retains the quantitative bound
\[
  q^{\mathrm{Alias}}_\gamma(M)
  \le
  \frac{2\overline\delta_{\mathcal P}(\pi)}{c(\gamma)}
\]
at finite regret.  This controls the evaluated mass of predictive aliases but
does not make $L(M(h))=\mathcal I^h$ well-defined everywhere.  Therefore no
exact monoid homomorphism or congruence interval is asserted at finite regret.
\end{remark}

This corollary is closer to internal containment than black-box transition
recovery, but its exact conclusion is intentionally limited to zero regret and
complete tests.  It is also relative to the evaluated history support unless
that support contains all admissible counterfactual histories.

\subsection{From recovery or selection to an internal quotient}
\label{subsec:controller-formalism}

There are therefore two possible upstream routes to a predictive state map.
Behavioural recovery in the style of Richens or Cifuentes must factor through
recurrent state.  Alternatively, complete zero-regret selection tests in the
style of Nayebi can force recurrent memory to refine predictive state on a
transition-closed tested support.  Once a surjective update-equivariant state
map exists, the remainder of the implication is the same:
\[
  \begin{aligned}
    \text{state-local predictive recovery}
    &\Longrightarrow Q\twoheadrightarrow\widehat X_{\mathcal I}\\
    &\Longrightarrow S_C\twoheadrightarrow S_{\mathcal I}\\
    &\Longrightarrow
    \Con(S_{\mathcal I})
    \cong[\ker\Phi_C,\top]\subseteq\Con(S_C).
  \end{aligned}
\]
The theorem below formalises the first route.
Corollary~\ref{cor:nayebi-refinement} formalises the second, subject to its
zero-regret, support, and completeness qualifications.

Fix a goal set $\mathcal G$.  Let a recurrent goal-conditioned controller be
\[
  C=(Q,q_0,d,\{\delta_{a,o}\}_{(a,o)\in A\times O},\pi).
\]
For a history $h=(a_1,o_1)\cdots(a_n,o_n)$, write
\[
  \delta_h
  =\delta_{a_n,o_n}\circ\cdots\circ\delta_{a_1,o_1},
  \qquad q_h=\delta_h(q_0).
\]
The non-dead state set is
\[
  Q^+=\{q_h:h\text{ admissible}\},
  \qquad Q=Q^+\mathbin{\dot\cup}\{d\}.
\]
The update maps are goal-independent and total, satisfy
$\delta_{a,o}(d)=d$, and send $q_h$ to $d$ whenever $h(a,o)$ is
impossible.  They are defined for every counterfactual intervention history,
not only histories generated by the controller's on-policy actions.  The
deterministic policy is a map $\pi:Q\times\mathcal G\to A$; if goal progress
affects recurrence, that memory must be included explicitly in $Q$.

For $q\in Q$, define its total continuation-policy signature by
\[
  B_C(q)(g,k)=\pi(\delta_k(q),g),
  \qquad
  g\in\mathcal G,\quad k\in(A\times O)^*.
\]
Using all counterfactual continuations, including those sent to $d$, makes
$B_C(q)$ a function of the controller state alone rather than of the history
chosen to represent it.

\begin{definition}[State-local behavioural recovery]
A recovery functional $\mathcal R$ is state-local when it is fixed uniformly
over the environment class under consideration and, for every admissible
history $h$,
\[
  \mathcal R(B_C(q_h))=\mathcal I^h.
\]
\end{definition}

This condition says that the policy information available from the recurrent
state determines the residual controlled process.  It is stronger than global
recovery from a policy oracle, because it requires the recoverable information
to factor through the controller state.

\begin{theorem}[Behaviour-to-state bridge]
\label{thm:controller-bridge}
Suppose a reachable recurrent goal-conditioned controller admits an exact
state-local recovery functional.  Then there is a surjection
\[
  L:Q\twoheadrightarrow\widehat X_{\mathcal I}
\]
satisfying
\[
  L\circ\delta_{a,o}=\tau_{a,o}\circ L
  \qquad((a,o)\in A\times O).
\]
\end{theorem}

\begin{proof}
For a state reached by an admissible history $h$, define
\[
  L(q_h)=\mathcal R(B_C(q_h))=\mathcal I^h,
\]
and send the controller's dead state to the predictive dead state.  This is
well-defined because equal controller states have the same complete
continuation-policy signature and hence the same recovered residual interface.
It is surjective because every element of $X_{\mathcal I}$ has an admissible
history representative, while the controller dead state is sent to
$\dagger$.

If $h(a,o)$ is admissible, state-local recovery at the successor gives
\[
  L(\delta_{a,o}q_h)
  =\mathcal I^{h(a,o)}
  =\tau_{a,o}(\mathcal I^h)
  =\tau_{a,o}L(q_h).
\]
For an impossible extension, both sides are the absorbing dead state.
At $d$ the equation holds because both $d$ and $\dagger$ are absorbing.
\end{proof}

Let
\[
  S_C=\langle\delta_{a,o}:(a,o)\in A\times O\rangle
\]
be the controller's update monoid.

\begin{corollary}[Controller congruence inheritance]
\label{cor:controller-con}
Under the hypotheses of Theorem~\ref{thm:controller-bridge}, there is a
surjective monoid homomorphism
\[
  \Phi_C:S_C\twoheadrightarrow S_{\mathcal I}.
\]
If $\kappa_C=\ker\Phi_C$, then
\[
  \Con(S_{\mathcal I})
  \cong[\kappa_C,\top]
  \subseteq\Con(S_C).
\]
Thus every $M_3$ or $N_5$ obstruction in the canonical controlled predictive
monoid is inherited by the recurrent controller.
\end{corollary}

\begin{proof}
For a word $w\in(A\times O)^*$, let $\delta_w$ and $\tau_w$ be the
corresponding composites and define
\[
  \Phi_C(\delta_w)=\tau_w.
\]
Iterating the intertwining equations gives
$L\circ\delta_w=\tau_w\circ L$.  If $\delta_w=\delta_v$, then
$\tau_w\circ L=\tau_v\circ L$, and surjectivity of $L$ gives
$\tau_w=\tau_v$.  Thus $\Phi_C$ is well-defined.  The common composition
convention makes it a homomorphism, and it is surjective because it hits every
generator $\tau_{a,o}$.  Congruence correspondence gives the displayed
interval isomorphism.
\end{proof}

\subsection{Approximation and the remaining gap}

The hypotheses above are deliberately exact.  Richens et al.'s finite-regret,
finite-depth theorem guarantees only approximate behavioural recovery, so it
does not by itself produce an exact congruence interval.  One extension is to
combine Cifuentes's common-support residual metric and predictive-state
separation theorem with monoid descent to study persistence of finite
$M_3/N_5$ witnesses under controlled approximation
\cite{CifuentesApprox2026}.

\section{Results and limitations}

The established theorem chain combines classical foundations with the new
monoid and congruence layer:
\begin{enumerate}[label=(\arabic*)]
  \item the conditional future laws of a process form a canonical predictive
        machine;
  \item every reachable exact predictive agent surjects onto that machine;
  \item the canonical transition monoid is a homomorphic image of the agent's
        transition monoid, so its congruence lattice is an interval in the
        agent's;
  \item projective transition dynamics omit observation likelihoods;
  \item a three-symbol Markov reset process already has an $M_3$ congruence
        interval despite admitting a transparent last-symbol predictor;
  \item fixed-task optimality does not force a predictive quotient;
  \item task-general policy recovery yields an internal predictive quotient only
        when the recovery map factors through a recurrent, counterfactually
        updateable controller state;
  \item in the zero-regret, complete-test limit, predictive memory refines
        predictive state; deterministic counterfactual updates then supply the
        exact quotient needed for descent.
\end{enumerate}

The contribution should therefore be read specifically at the monoid and
lattice level: monoid descent, the inherited congruence interval, its
$M_3/N_5$ witnesses, and the reset-monoid limitation.  The causal-state,
$\varepsilon$-transducer, Hankel,
behavioural-recovery, and predictive-memory results used to reach that layer are
prior work.  Failure of distributivity can be a process-forced property of every
exact predictor's transition monoid, but without an additional criterion
separating reset effects from stronger algebraic constraints it is not a
certificate of feature entanglement or interpretability failure.

\begin{thebibliography}{99}

\bibitem{Fliess1974}
M.~Fliess,
\emph{Matrices de Hankel.}
Journal de Math\'ematiques Pures et Appliqu\'ees 53 (1974), 197--222.

\bibitem{Jaeger2000}
H.~Jaeger,
\emph{Observable operator models for discrete stochastic time series.}
Neural Computation 12 (2000), 1371--1398.

\bibitem{CrutchfieldYoung1989}
J.~P.~Crutchfield and K.~Young,
\emph{Inferring statistical complexity.}
Physical Review Letters 63 (1989), 105--108.

\bibitem{ShaliziCrutchfield2001}
C.~R.~Shalizi and J.~P.~Crutchfield,
\emph{Computational mechanics: Pattern and prediction, structure and simplicity.}
Journal of Statistical Physics 104 (2001), 817--879.

\bibitem{BarnettCrutchfield2015}
N.~Barnett and J.~P.~Crutchfield,
\emph{Computational mechanics of input--output processes: Structured
transformations and the $\varepsilon$-transducer.}
Journal of Statistical Physics 161 (2015), 404--451.

\bibitem{DaveyPriestley2002}
B.~A.~Davey and H.~A.~Priestley,
\emph{Introduction to Lattices and Order}, 2nd ed.
Cambridge University Press, 2002.

\bibitem{Howie1995}
J.~M.~Howie,
\emph{Fundamentals of Semigroup Theory.}
Oxford University Press, 1995.

\bibitem{Richens2025}
J.~Richens, D.~Abel, A.~Bellot, and T.~Everitt,
\emph{General agents contain world models.}
arXiv:2506.01622, 2025.

\bibitem{CifuentesApprox2026}
S.~Cifuentes,
\emph{Approximate homomorphism of transducers and minimal implementations.}
Unpublished draft, June 2026.

\bibitem{Cifuentes2026}
S.~Cifuentes,
\emph{General agents contain world models, even under partial observability and
stochasticity.}
arXiv:2602.03146, 2026.

\bibitem{Nayebi2026}
A.~Nayebi,
\emph{What capable agents must know: Selection theorems for robust
decision-making under uncertainty.}
arXiv:2603.02491, 2026.

\bibitem{LuEtAl2026}
Y.~Lu, Y.~Wu, X.~Lu, and T.~Li,
\emph{World models in pieces: Structural certification for general agents.}
arXiv:2606.24842, 2026.

\bibitem{ConantAshby1970}
R.~C.~Conant and W.~R.~Ashby,
\emph{Every good regulator of a system must be a model of that system.}
International Journal of Systems Science 1 (1970), 89--97.

\end{thebibliography}

\end{document}
